\documentclass{article}
\input{../../../../lib.sty}

\title{\texttt{fn match\_num\_groups\_predicate}}
\author{\MichaelShoemateAuthor}
\begin{document}
\maketitle


\contrib

Proves soundness of \rustdoc{transformations/make_stable_lazyframe/truncate/matching/fn}{match\_num\_groups\_predicate}
in \asOfCommit{mod.rs}{f5bb719}.

\section{Hoare Triple}
\subsection*{Precondition}
\subsubsection*{Compiler-verified}
Types matching pseudocode.

\subsubsection*{Caller-verified}
None

\subsection*{Pseudocode}
\label{sec:python-pseudocode}
\lstinputlisting[language=Python,firstline=2,escapechar=|]{./pseudocode/match_num_groups_predicate.py}

\subsection*{Postcondition}
\begin{theorem}[Postcondition]
    If \texttt{ranks} is a dense ranking of grouping columns,
    and \texttt{partition\_by} is a singleton of \texttt{identifier},
    then returns the bound on per-identifier contributions,
    or an error if the truncation is mis-specified.
\end{theorem}

\begin{proof}
    Line \ref{line:check_rank} first distinguishes a dense-rank truncation
    from other expressions. Non-ranking expressions are not treated as
    number-of-groups truncations; malformed rank expressions are rejected
    by the subsequent checks.

    Line \ref{line:check_identifier} checks that \texttt{partition\_by} is a singleton of the \texttt{identifier}, meeting the conditions of the postcondition.

    We now check whether the truncation predicate is well-defined, on lines \ref{line:check_method} and \ref{line:check_input}.

    Finally, line \ref{line:extract_grouping_columns} extracts the grouping columns from the predicate.
    Special consideration is made for structs, which are considered multiple grouping columns.

    Further consideration is made if the first field is a hash of the second field, on line \ref{line:check_hash}.
    Field aliases are ignored when checking this relationship: they prevent duplicate
    field names in Polars without changing the values.
    The shape check on line \ref{line:check_hash_shape} requires exactly two fields
    (the hash and the complete grouping key), and exactly one struct-valued hash input. Without these checks, unrecognized
    fields could affect the dense ranks but be omitted from the returned bound.
    Before the bound is used, the caller restricts each dense-rank grouping
    key to supported scalar dtypes, excluding nested and categorical keys
    in both hashed and non-hashed ranks.
    Physical categorical codes can depend on input row order, so ranking or hashing
    these codes need not be permutation invariant.
    This deterministically orders grouping columns by hash and then by the original keys; it does not produce a uniformly random sample.
    The second field of grouping columns is kept in the ranker for the rare case that there are hash collisions,
    which prevents different combinations of grouping keys from being assigned the same rank,
    preventing more distinct groups being kept than is permitted.

    This predicate corresponds to a \texttt{num\_groups} truncation predicate,
    because the over expression groups by the \texttt{identifier} column,
    and within each group, a dense ranking is applied to unique combinations of the grouping columns.
    If the only rows kept are those corresponding to grouping keys assigned dense ranks at most \texttt{threshold},
    then each user identifier will have at most \texttt{threshold} unique combinations of the grouping columns
    after filtering by the predicate.

    Therefore the bound on user contributions constructed on line \ref{line:bound} is valid.
\end{proof}
\end{document}
